\documentclass[11pt,a4paper]{letter}
\usepackage[utf8]{inputenc}
\usepackage[T1]{fontenc}
\usepackage{geometry}
\usepackage{enumitem}

% Ustawienia marginesów pozwalające zmieścić więcej treści na 1 stronie
\geometry{
  top=2cm,
  bottom=2cm,
  left=2.5cm,
  right=2.5cm
}

% Dane nadawcy
\address{
    \textbf{Maciej Tadej}\\
    Institute of Mathematics\\
    University of Wrocław\\
    pl. Grunwaldzki 2\\
    50-384 Wrocław, Poland\\
    Email: maciej.tadej@math.uni.wroc.pl
}

\date{\today}

\begin{document}

\begin{letter}{
    The Editorial Board of the journal: 
    \textit{ \\ Communications in Nonlinear Science \\ and Numerical Simulation, Elsevier}
}

\opening{Dear Editor,}

We are pleased to submit the manuscript entitled \textbf{``Extinction and persistence criteria in non-local Klausmeier models of vegetation dynamics on flat landscapes''} for consideration for publication in \textit{Communications in Nonlinear Science and Numerical Simulation}.

The manuscript investigates how non-local dispersal qualitatively alters extinction–persistence transitions, bifurcation structure, and spatial morphology in a nonlinear Klausmeier model of vegetation pattern formation. 

We derive a spectral persistence criterion that defines a critical habitat size and demonstrate that this threshold depends sensitively on the shape of the dispersal kernel. Further, we identify a critical minimum biomass threshold that guarantees extinction irrespective of habitat size.  These analytical results are complemented by numerical bifurcation analysis, revealing complex solution branches, kernel-dependent resilience, and patterned states persisting below the classical kinetic rainfall threshold. Together, the results highlight non-locality as a genuine nonlinear mechanism that enhances persistence and reshapes spatial organization in finite systems.

While the manuscript includes rigorous mathematical analysis to support the results, this analysis is used primarily to explain observed nonlinear phenomena and stability transitions. We believe the combination of nonlinear dynamics, bifurcation structure, and physically motivated non-local effects aligns well with the scope and readership of Physica D.

The manuscript has not been published previously and is not under consideration elsewhere. All authors have approved the submission.

Thank you for your consideration.

\flushright{Sincerely, \\ Maciej Tadej \\ Ricardo Martinez-Garc\'ia \\ Michael Hecht}

\end{letter}
\end{document}